% ===========================================================================
\chapter{Introdução e Objetivo do Projeto}
\label{cap_introducao_objetivo}
% ===========================================================================

Este relatório consolida a busca sistemática de anterioridade e prospecção tecnológica desenvolvida para a concepção, proteção intelectual e validação do sistema \textbf{SIGEA-GTT} (\textit{Sistema Inteligente de Gestão de Eventos Adversos -- Global Trigger Tool}). A solução foi projetada como um sistema de monitoramento, rastreamento e auditoria de riscos clínicos hospitalares com aplicação educacional síncrona voltada à promoção da segurança do paciente nas instituições de saúde \cite{ihi2009}.

O escopo central do projeto reside no desenvolvimento de uma arquitetura de software para a identificação ativa e o rastreamento retrospectivo de eventos adversos (EAs) e incidentes clínicos a partir da metodologia internacional de gatilhos globais (\textit{Global Trigger Tool} -- GTT), padronizada pelo \textit{Institute for Healthcare Improvement} (IHI). 

% ===========================================================================
\section{Diferencial Inovador: Abordagem Não-Punitiva}
\label{sec_diferencial_inovador}
% ===========================================================================

O principal diferencial inventivo da plataforma repousa na sua \textbf{aplicação pedagógica}: consolidar-se como um ambiente computacional simulado, interativo e estritamente não-punitivo para a formação clínica e tomada de decisão de graduandos em Enfermagem e Medicina.

Historicamente, as ferramentas de auditoria e monitoramento de eventos adversos foram concebidas com foco exclusivo em ambientes corporativos hospitalares, voltadas à apuração de responsabilidades administrativas, controle financeiro de glosas hospitalares e auditoria de faturamento de operadoras de planos de saúde. Essa abordagem tradicional fomenta uma cultura de medo e omissão entre profissionais e estudantes, inibindo o reporte de falhas e ocultando a gênese sistêmica dos incidentes assistenciais.

Ao estruturar a leitura sistemática de prontuários, a indexação determinística dos 53 gatilhos clínicos canônicos do IHI e a emissão de painéis gerenciais e relatórios sintético-analíticos em um ambiente controlado (\textit{sandbox}), o SIGEA-GTT promove a transição paradigmática de uma cultura hospitalar punitiva para uma abordagem sistemática de mitigação de riscos, análise de causa raiz e melhoria contínua da qualidade do cuidado.

% ===========================================================================
\section{Delimitação Institucional e Ética}
\label{sec_delimitacao_etica}
% ===========================================================================

Em conformidade com a Lei Geral de Proteção de Dados Pessoais (Lei n.º 13.709/2018 -- LGPD) \cite{lgpd2018}, o SIGEA-GTT opera exclusivamente sobre bases sintéticas e fictícias de prontuários de alta fidelidade persistidos em colunas semiestruturadas \texttt{JSONB}. A plataforma mantém isolamento físico e lógico absoluto em relação às redes assistenciais e bases de dados reais de pacientes do Hospital Universitário da Universidade Federal de Sergipe (HU-UFS) ou do Sistema Único de Saúde (SUS), garantindo a conformidade ética e a segurança jurídica necessárias para a pesquisa e inovação tecnológica universitária.
